\begin{Example}
\label{ex:fourier-heaviside-regularization}
不能以经典积分形式定义的 Fourier 变换, 常常可以作为正则化的极限求出. 例如, 考虑 Heaviside 函数 $H(x)$. 显然, 不能将 Fourier 变换定义为
\begin{equation}\label{eq:heaviside-nonconvergent-fourier-integral}
\dfrac1{\sqrt{2\pi}}\int_0^\infty e^{-i\xi x}\,\mathrm{d}x.
\end{equation}
然而, 注意到
\begin{equation}\label{eq:heaviside-exponential-regularization}
H(x)=\lim_{\epsilon\to0^+}H(x)e^{-\epsilon x}
\end{equation}
在缓增分布意义下成立, 因而
\begin{equation}\label{eq:heaviside-fourier-regularized-limit}
\mathcal{F}[H](\xi)=\dfrac1{\sqrt{2\pi}}\lim_{\epsilon\to0^+}\int_0^\infty e^{-\epsilon x-i\xi x}\,\mathrm{d}x=\lim_{\epsilon\to0^+}\dfrac1{\sqrt{2\pi}}\dfrac1{\epsilon+i\xi}.
\end{equation}
将这一极限作用于试验函数, 可以更明确地求出它. 对任意 $\delta>0$, 有
\begin{equation}\label{eq:heaviside-fourier-principal-value-limit}
\begin{aligned}
\lim_{\epsilon\to0^+}\int_{-\infty}^\infty\dfrac{\phi(\xi)}{\epsilon+i\xi}\,\mathrm{d}\xi
&=\int_{|\xi|>\delta}\dfrac{\phi(\xi)}{i\xi}\,\mathrm{d}\xi
+\int_{-\delta}^{\delta}\dfrac{\phi(\xi)-\phi(0)}{i\xi}\,\mathrm{d}\xi\\
&\quad+\lim_{\epsilon\to0^+}\int_{-\delta}^{\delta}\dfrac{\phi(0)}{\epsilon+i\xi}\,\mathrm{d}\xi\\
&=\int_{|\xi|>\delta}\dfrac{\phi(\xi)}{i\xi}\,\mathrm{d}\xi
+\int_{-\delta}^{\delta}\dfrac{\phi(\xi)-\phi(0)}{i\xi}\,\mathrm{d}\xi+\pi\phi(0).
\end{aligned}
\end{equation}
得到
\begin{equation}\label{eq:heaviside-fourier-principal-value-formula}
\mathcal{F}[H]=\dfrac1{\sqrt{2\pi}}\left(\dfrac1{i\xi}+\pi\delta\right),
\end{equation}
其中 $1/(i\xi)$ 理解为主值.
\end{Example}

\begin{Example}
\label{ex:fourier-polynomial-growth-truncation}
设 $f$ 为任意在无穷远处具有多项式增长的连续函数. 那么在缓增分布意义下, $f$ 为下列函数当 $M\to\infty$ 时的极限
\begin{equation}\label{eq:polynomial-growth-function-truncation}
f_M(\bm{x})=\begin{cases}f(\bm{x}),&|\bm{x}|\leq M\\0,&|\bm{x}|>M.\end{cases}
\end{equation}
由此得到, 在缓增分布意义下,
\begin{equation}\label{eq:polynomial-growth-fourier-truncation-limit}
\widehat f(\bm{\xi})=(2\pi)^{-m/2}\lim_{M\to\infty}\int_{|\bm{x}|\leq M}f(\bm{x})e^{-i\bm{\xi}\cdot\bm{x}}\,\mathrm{d}\bm{x}.
\end{equation}
特别地, 若 $f$ 在无穷远处可积, 则 $f$ 作为分布的 Fourier 变换与普通 Fourier 变换一致. 因此, 计算具有多项式增长的函数的 Fourier 变换的另一种方法, 是用 $f(\bm{x})\exp(-\epsilon|\bm{x}|^2)$ 这样的可积函数逼近它们. 例子见 Problem~\ref{ex:fourier-polynomial-growth-regularizations}.
\end{Example}

\begin{Example}
\label{ex:fourier-uniform-sphere-mass}
设 $\delta(r-a)$ 表示半径为 $a$ 的球面上的均匀质量分布, 即
\begin{equation}\label{eq:uniform-sphere-mass-pairing}
(\delta(r-a),\phi)=\int_{|\bm{x}|=a}\phi(\bm{x})\,\mathrm{d}S.
\end{equation}
(当然, 这与前面将 "$\delta$" 用作 $\mathbb{R}^m$ 上分布的方式不一致, 但这是一种标准的记号滥用, 读者应当逐渐习惯.) 那么 $\delta(r-a)$ 的 Fourier 变换由 \eqref{eq:compact-distribution-fourier-entire-formula} 给出
\begin{equation}\label{eq:uniform-sphere-mass-fourier-integral}
\mathcal{F}[\delta(r-a)](\bm{\xi})=(2\pi)^{-m/2}\int_{|\bm{x}|=a}e^{-i\bm{\xi}\cdot\bm{x}}\,\mathrm{d}S.
\end{equation}
要在 $m=3$ 时计算这一表达式. 使用以 $\bm{\xi}$ 的方向为轴向的极坐标, 使 $\bm{\xi}\cdot\bm{x}=a|\bm{\xi}|\cos\theta$; 用 $\rho$ 表示 $|\bm{\xi}|$. 于是得到
\begin{equation}\label{eq:uniform-sphere-mass-three-dimensional-fourier}
\begin{aligned}
\mathcal{F}[\delta(r-a)](\bm{\xi})
&=(2\pi)^{-3/2}a^2\int_0^\pi\int_0^{2\pi}e^{-ia\rho\cos\theta}\sin\theta\,\mathrm{d}\phi\,\mathrm{d}\theta\\
&=\sqrt{\dfrac2\pi}a\dfrac{\sin a\rho}\rho.
\end{aligned}
\end{equation}

\end{Example}

\begin{Example}
\label{ex:fourier-direct-product}
直积的 Fourier 变换是各 Fourier 变换的直积. 为说明这一点, 只需对一个稠密的试验函数集合证明相等. 有
\begin{equation}\label{eq:fourier-direct-product-separable-proof}
\begin{aligned}
(\widehat f(\bm{\xi})\widehat g(\bm{\eta}),\widehat\phi(\bm{\xi})\widehat\psi(\bm{\eta}))
&=(\widehat f,\widehat\phi)(\widehat g,\widehat\psi)\\
&=(f,\phi)(g,\psi)\\
&=(f(\bm{x})g(\bm{y}),\phi(\bm{x})\psi(\bm{y})).
\end{aligned}
\end{equation}

\end{Example}
